\documentclass[a4paper,11pt]{article}
\usepackage[utf8]{inputenc}
\usepackage[bahasa]{babel}
\usepackage{geometry}
\usepackage{array}
\usepackage{longtable}
\usepackage{booktabs}

% Pengaturan Margin
\geometry{
    top=2.5cm,
    bottom=2.5cm,
    left=3cm,
    right=2.5cm
}

% Definisi kolom tabel
\newcolumntype{L}[1]{>{\raggedright\let\newline\\\arraybackslash\hspace{0pt}}m{#1}}
\newcolumntype{C}[1]{>{\centering\let\newline\\\arraybackslash\hspace{0pt}}m{#1}}

\begin{document}

% --- KOP SURAT ---
\begin{center}
    \textbf{\large INSTITUT TEKNOLOGI BANDUNG}\\
    \textbf{\large SEKOLAH TEKNIK ELEKTRO DAN INFORMATIKA}\\
    \textbf{PROGRAM STUDI MAGISTER INFORMATIKA}\\
    \vspace{0.2cm}
    \hrule height 1pt
    \vspace{0.1cm}
    \hrule height 0.5pt
    \vspace{0.5cm}
    \textbf{\large JADWAL KEGIATAN PENELITIAN TESIS}\\
    \textbf{Periode: 50 Minggu}
\end{center}
\vspace{0.5cm}

% --- DATA MAHASISWA ---
\noindent
\begin{tabular}{L{4cm} L{0.5cm} L{10cm}}
    \textbf{Nama} & : & Jatniko Nur Mutaqin \\
    \textbf{NIM} & : & 23523314 \\
    \textbf{Judul Tesis} & : & Restorasi Dokumen Terdegradasi Menggunakan Generative Adversarial Network dengan Diskriminator Dual-Modal dan Optimasi Loss Function Berorientasi HTR
\end{tabular}

\vspace{0.5cm}

% --- TABEL JADWAL KEGIATAN ---
\begin{longtable}{|C{2.5cm}|L{5cm}|L{7cm}|}
\hline
\textbf{Minggu Ke} & \textbf{Kegiatan} & \textbf{Keterangan} \\
\hline
\endfirsthead
\hline
\textbf{Minggu Ke} & \textbf{Kegiatan} & \textbf{Keterangan} \\
\hline
\endhead

1 -- 4 & Studi Literatur & Kajian pustaka mengenai GAN, restorasi dokumen, dan \textit{Handwritten Text Recognition} (HTR) \\
\hline

5 -- 8 & Analisis Kebutuhan & Identifikasi permasalahan, perumusan hipotesis, dan penentuan ruang lingkup penelitian \\
\hline

9 -- 12 & Perancangan Arsitektur & Desain arsitektur generator U-Net dan diskriminator dual-modal \\
\hline

13 -- 16 & Pengembangan Dataset & Pengumpulan dan pembuatan dataset dokumen sintetis serta \textit{ground truth} \\
\hline

17 -- 20 & Implementasi Model & Pengembangan kode program generator dan diskriminator menggunakan PyTorch \\
\hline

21 -- 24 & Implementasi \textit{Loss Function} & Integrasi lima komponen fungsi kerugian multiobjektif \\
\hline

25 -- 28 & Eksperimen Awal & Pelatihan model baseline dan validasi fungsionalitas sistem \\
\hline

29 -- 32 & Optimasi \textit{Hyperparameter} & \textit{Grid search} dan \textit{tuning} parameter pelatihan model \\
\hline

33 -- 36 & Studi Ablasi & Eksperimen ablasi untuk validasi kontribusi setiap komponen arsitektur \\
\hline

37 -- 40 & Evaluasi Kuantitatif & Pengukuran metrik PSNR, SSIM, CER, dan analisis statistik \\
\hline

41 -- 44 & Analisis Hasil & Interpretasi hasil eksperimen dan validasi hipotesis penelitian \\
\hline

45 -- 48 & Penulisan Tesis & Penyusunan dokumen tesis Bab I hingga VI \\
\hline

49 -- 50 & Finalisasi \& Publikasi & Revisi akhir, persiapan seminar, dan submisi makalah ke jurnal \\
\hline

\end{longtable}

\vspace{1cm}

% --- TANDA TANGAN ---
\begin{table}[h!]
    \centering
    \begin{tabular}{p{7cm} p{7cm}}
        & Bandung, 12 Desember 2025 \\
        Menyetujui, & Mahasiswa, \\
        Pembimbing & \\
        \vspace{2.5cm} & \vspace{2.5cm} \\
        \textbf{Ir. I Gusti Bagus Baskara Nugraha, ST, MT, Ph.D} & \textbf{Jatniko Nur Mutaqin} \\
        NIP. ........................... & NIM. 23523314
    \end{tabular}
\end{table}

\end{document}
